%! Characteristics of Logarithmic Functions — Unit 7, Lesson 7.0 — Warm-Up
\documentclass[10pt]{article}
\usepackage{algebra2-article}
\usepackage{algebra2-boxes}
\newcommand{\TallMath}[1]{$\displaystyle #1\rule[-1.4em]{0pt}{3.2em}$}
\newcommand{\lgrid}[4]{%
  \draw[step=1, linegray!40, very thin] (#1,#3) grid (#2,#4);
  \draw[->, semithick, charcoal] (#1,0)--(#2,0) node[below right, font=\scriptsize]{$x$};
  \draw[->, semithick, charcoal] (0,#3)--(0,#4) node[left, font=\scriptsize]{$y$};}

\begin{document}

\pageheader{Unit 7, Lesson 7.0}{Warm-Up}
\namedateperiod

\vspace{0.08in}

\begin{notesbox}{Getting Ready: Skills We'll Use Today}
\small These three skills power today's lesson on reading logarithmic graphs. Work quickly.

\vspace{3pt}
\textbf{1. Exponents --- forwards, then backwards.} Evaluate each one.
\begin{center}
\renewcommand{\arraystretch}{1.4}
\begin{tabular}{@{}l l l@{}}
(a) \; $2^3=$ \blank{1.6cm} & (b) \; $2^0=$ \blank{1.6cm} & (c) \; $2^{-2}=$ \blank{1.6cm} \\
\end{tabular}
\end{center}
Now run it \emph{backwards}. Fill in the missing exponent.
\begin{center}
\renewcommand{\arraystretch}{1.4}
\begin{tabular}{@{}l l l@{}}
(d) \; $2^{\square}=32$, \; $\square=$ \blank{1.4cm} & (e) \; $2^{\square}=\tfrac18$, \; $\square=$ \blank{1.4cm}
 & (f) \; $2^{\square}=-8$, \; $\square=$ \blank{1.4cm} \\
\end{tabular}
\end{center}
One of (d)--(f) is impossible. Which one, and what is it about \emph{powers of 2} that makes it
impossible? \writeline

\vspace{3pt}
\textbf{2. Read the exponential parent.} The graph is $y=2^x$ from Unit~6.
\begin{center}
\begin{minipage}[c]{0.40\linewidth}
\centering
\begin{tikzpicture}[scale=0.42]
  \lgrid{-4}{3}{-1}{5}
  \draw[navy, dashed, semithick] (-4,0)--(3,0);
  \begin{scope}
    \clip (-4,-1) rectangle (3,5);
    \draw[very thick, forest, domain=-4:2.32, samples=90, smooth] plot (\x,{pow(2,\x)});
  \end{scope}
  \fill[charcoal] (0,1) circle (2.6pt);
  \fill[charcoal] (1,2) circle (2.6pt);
  \fill[charcoal] (2,4) circle (2.6pt);
\end{tikzpicture}
\end{minipage}
\begin{minipage}[c]{0.55\linewidth}
\small
Domain: \blank{2.6cm} \\[4pt]
Range: \blank{2.6cm} \\[4pt]
Horizontal asymptote: $y=$ \blank{1.4cm} \\[4pt]
$y$-intercept: \blank{1.8cm} \quad $x$-intercept: \blank{1.8cm}
\end{minipage}
\end{center}

\vspace{2pt}
\textbf{3. Two tables that are hiding something.} Complete the first table. The second table is the
first one \emph{with its rows traded}.
\begin{center}
\renewcommand{\arraystretch}{1.25}
\begin{tabular}{|c|c|c|c|c|c|}
\hline
\multicolumn{6}{|c|}{\small $y=2^x$} \\ \hline
$x$ & $-2$ & $0$ & $1$ & $2$ & $3$ \\ \hline
$y$ & \blank{0.8cm} & \blank{0.8cm} & \blank{0.8cm} & \blank{0.8cm} & \blank{0.8cm} \\ \hline
\end{tabular}
\hspace{0.15in}
\begin{tabular}{|c|c|c|c|c|c|}
\hline
\multicolumn{6}{|c|}{\small the same table, rows traded} \\ \hline
$x$ & \blank{0.8cm} & \blank{0.8cm} & \blank{0.8cm} & \blank{0.8cm} & \blank{0.8cm} \\ \hline
$y$ & $-2$ & $0$ & $1$ & $2$ & $3$ \\ \hline
\end{tabular}
\end{center}
Look only at the \emph{inputs} of the second table. What do they all have in common? \writeline

\end{notesbox}

\end{document}
